%% Section 4: Topologies and modulation
\section{Topologies and Modulation}
\label{sec:topologies}

\subsection{H4 inverter}
The H4 inverter comprises four switches $S_1$--$S_4$ with anti-parallel diodes (Fig.~\ref{fig:schematic} without $S_5$, $S_6$). The pairs $(S_1,S_4)$ and $(S_2,S_3)$ are the diagonal pairs that deliver positive and negative output voltage, respectively.

\emph{Bipolar SPWM:} a sinusoidal reference is compared with a triangular carrier; the diagonal pairs are switched complementarily at $\fsw$. The output alternates between $+\Vdc$ and $-\Vdc$ (two levels) and, ideally, $\vCM=\Vdc/2$ at all times.

\emph{Unipolar SPWM:} the legs are modulated with sinusoidal references of opposite sign. The output has three levels ($+\Vdc$, $0$, $-\Vdc$), which halves the effective ripple frequency seen by the filter and improves efficiency, but the freewheeling states connect both poles to the same rail so that $\vCM$ steps between $0$, $\Vdc/2$ and $\Vdc$ at the carrier frequency (Table~\ref{tab:cmv}).

\subsection{HERIC inverter}
The HERIC inverter adds two switches $S_5$, $S_6$ with diodes between the bridge output nodes A and B on the AC side \cite{Schmidt2003}. During the active intervals the bridge delivers $\pm\Vdc$ as in the H4 inverter; during the zero-voltage intervals all bridge switches are off and the filter current freewheels through the AC bypass ($S_5$ for one polarity, $S_6$ for the other), which disconnects the AC side from the PV array. Because the DC side then floats, the CMV remains, ideally, at $\Vdc/2$ in all four operating states, which is the origin of the low-leakage property \cite{Araujo2010,Sundar2016}. In real circuits the decoupling is imperfect owing to the parasitic capacitances of the switches and to dead time, so a residual CMV component remains.

\begin{table}[!t]
\renewcommand{\arraystretch}{1.2}
\caption{Ideal CMV behavior of the configurations studied}
\label{tab:cmv}
\centering
\footnotesize
\begin{adjustbox}{max width=\columnwidth}
\begin{tabular}{p{3.0cm}p{2.8cm}p{2.4cm}p{3.0cm}}
\toprule
Configuration & Output levels & Ideal $\vCM$ & Carrier-frequency content of $\vCM$ \\
\midrule
H4, bipolar SPWM & 2 ($\pm\Vdc$) & $\Vdc/2$ & None (ideal) \\
H4, unipolar SPWM & 3 ($\pm\Vdc$, 0) & $0$, $\Vdc/2$, $\Vdc$ & Yes \\
HERIC, unipolar SPWM & 3 ($\pm\Vdc$, 0) & $\Vdc/2$ & None (ideal) \\
HERIC, bipolar SPWM & 2 ($\pm\Vdc$) with bypass & $\Vdc/2$ & None (ideal) \\
\bottomrule
\end{tabular}
\end{adjustbox}
\end{table}

The ``HERIC bipolar'' configuration studied here applies bipolar SPWM to the bridge while the bypass switches are driven at line frequency; it is not the conventional HERIC modulation, which uses unipolar PWM, and it is included as a variant because it was found to give the lowest leakage in the simulations.

\subsection{Modulation and ripple considerations}
In SPWM the switching signals are generated by comparing a sinusoidal reference $v_r=m_a\sin(\omega t)$, $0<m_a\le1$, with a triangular carrier of frequency $\fsw$ and unit amplitude. For bipolar SPWM a single reference drives the diagonal pairs complementarily; for unipolar SPWM the references $+v_r$ and $-v_r$ drive the two legs separately, so that the bridge output has three levels and its spectrum is centered on $2\fsw$ rather than $\fsw$. For a single bridge leg the duty cycle is $D(\omega t)=m_a\sin(\omega t)$ and the inductor current ripple is
\begin{equation}
\Delta I_{pp}=\frac{D\,T_s\,(\Vdc-V_{o})}{L},\qquad V_o=\Vdc D,
\label{eq:ripple}
\end{equation}
whose maximum with respect to the duty cycle, at $D=1/2$, gives $\Delta I_{pp,\max}=\Vdc/(4\fsw L)$. This bound is useful in this study because the leakage-current sweep spans a ratio of 40 in $\fsw$: at the lowest frequencies the ripple, and hence the low-order THD, is large (Section~\ref{sec:results}). The design rules for LCL filters in \cite{Reznik2014} (total inductance limited to about 10\% of the base impedance, shunt capacitor absorbing about 5\% of the rated reactive power, and the resonance placed between $10f_g$ and $0.5\fsw$) are reproduced here only for context, since the L filter used in this study has no shunt branch.

\subsection{Choice of output filter}
An L filter was adopted. Shunt capacitors in LC and LCL filters reduce the harmonic content more effectively, and \cite{Reznik2014} provides the corresponding design rules (e.g., a resonance between ten times the grid frequency and half the switching frequency). However, a shunt capacitor referred to the grounded neutral adds a high-frequency CM path, and it complicates the CM model of Section~\ref{sec:cm}. The L filter keeps the CM impedance determined by $\Leq$ alone and thus isolates the effect of the switching frequency. The inductor ripple at the worst case follows the standard result $\Delta I_{pp,\max}=\Vdc/(4\fsw L)$, from which $L$ can be chosen for a given ripple limit at the lowest swept frequency. A comparison with LC and LCL filters is left for future work.
